\documentclass[12pt,a4paper]{exam}
\usepackage[utf8]{inputenc}
\usepackage{amsmath,amssymb,amsthm}
\usepackage[margin=1in]{geometry}
\usepackage{graphicx}
\usepackage[colorlinks=true]{hyperref}
\pagestyle{headandfoot}
\firstpageheader{MIT OpenCourseWare}{}{\textbf{EXTREMELY HARD -- MIT LEVEL}}
\runningheader{MIT OpenCourseWare}{}{\textbf{EXTREMELY HARD -- MIT LEVEL}}
\firstpagefooter{}{Page \thepage}{}
\runningfooter{}{Page \thepage}{}

\begin{document}

\begin{center}
{\LARGE \textbf{MIT OpenCourseWare}}\\2.005\\Thermal-Fluids Engineering I
\vspace{0.5cm}
{2.005} --- {Thermal-Fluids Engineering I}
\vspace{0.3cm}
May 2026
\vspace{0.3cm}
\textbf{EXTREMELY HARD -- MIT LEVEL}
\end{center}

\begin{center}
\textbf{Time Limit: 3 hours}
\end{center}

\begin{center}
\textbf{Total Marks: 100}
\end{center}

\vspace{0.5cm}

\begin{questions}

\question[3]
Construct an original proof-level problem involving Thermodynamics for 2.005 (Thermal-Fluids Engineering I) and provide a complete solution. Your problem should be at the level of a graduate qualifying exam.

\vspace{0.3cm}

\question[4]
Analyze the limitations of standard approaches to Fluid Mechanics when applied to edge cases in 2.005. Construct explicit counterexamples showing where intuitive reasoning fails.

\vspace{0.3cm}

\question[4]
Prove the fundamental theorem concerning Heat Transfer in the context of 2.005 (Thermal-Fluids Engineering I). State the theorem precisely, provide a complete proof, and discuss three important corollaries.

\vspace{0.3cm}

\question[5]
Prove the fundamental theorem concerning Boundary Layers in the context of 2.005 (Thermal-Fluids Engineering I). State the theorem precisely, provide a complete proof, and discuss three important corollaries.

\vspace{0.3cm}

\question[5]
Prove the fundamental theorem concerning Turbulence in the context of 2.005 (Thermal-Fluids Engineering I). State the theorem precisely, provide a complete proof, and discuss three important corollaries.

\vspace{0.3cm}

\question[2]
Prove the fundamental theorem concerning Thermodynamics in the context of 2.005 (Thermal-Fluids Engineering I). State the theorem precisely, provide a complete proof, and discuss three important corollaries.

\vspace{0.3cm}

\question[4]
Construct an original proof-level problem involving Fluid Mechanics for 2.005 (Thermal-Fluids Engineering I) and provide a complete solution. Your problem should be at the level of a graduate qualifying exam.

\vspace{0.3cm}

\question[6]
Analyze the limitations of standard approaches to Heat Transfer when applied to edge cases in 2.005. Construct explicit counterexamples showing where intuitive reasoning fails.

\vspace{0.3cm}

\question[5]
Analyze the limitations of standard approaches to Boundary Layers when applied to edge cases in 2.005. Construct explicit counterexamples showing where intuitive reasoning fails.

\vspace{0.3cm}

\question[4]
Analyze the limitations of standard approaches to Turbulence when applied to edge cases in 2.005. Construct explicit counterexamples showing where intuitive reasoning fails.

\vspace{0.3cm}

\question[4]
Prove the fundamental theorem concerning Thermodynamics in the context of 2.005 (Thermal-Fluids Engineering I). State the theorem precisely, provide a complete proof, and discuss three important corollaries.

\vspace{0.3cm}

\question[6]
Construct an original proof-level problem involving Fluid Mechanics for 2.005 (Thermal-Fluids Engineering I) and provide a complete solution. Your problem should be at the level of a graduate qualifying exam.

\vspace{0.3cm}

\question[4]
Prove the fundamental theorem concerning Heat Transfer in the context of 2.005 (Thermal-Fluids Engineering I). State the theorem precisely, provide a complete proof, and discuss three important corollaries.

\vspace{0.3cm}

\question[4]
Prove the fundamental theorem concerning Boundary Layers in the context of 2.005 (Thermal-Fluids Engineering I). State the theorem precisely, provide a complete proof, and discuss three important corollaries.

\vspace{0.3cm}

\question[4]
Prove the fundamental theorem concerning Turbulence in the context of 2.005 (Thermal-Fluids Engineering I). State the theorem precisely, provide a complete proof, and discuss three important corollaries.

\vspace{0.3cm}

\question[4]
Construct an original proof-level problem involving Thermodynamics for 2.005 (Thermal-Fluids Engineering I) and provide a complete solution. Your problem should be at the level of a graduate qualifying exam.

\vspace{0.3cm}

\question[3]
Analyze the limitations of standard approaches to Fluid Mechanics when applied to edge cases in 2.005. Construct explicit counterexamples showing where intuitive reasoning fails.

\vspace{0.3cm}

\question[3]
Construct an original proof-level problem involving Heat Transfer for 2.005 (Thermal-Fluids Engineering I) and provide a complete solution. Your problem should be at the level of a graduate qualifying exam.

\vspace{0.3cm}

\question[3]
Analyze the limitations of standard approaches to Boundary Layers when applied to edge cases in 2.005. Construct explicit counterexamples showing where intuitive reasoning fails.

\vspace{0.3cm}

\question[5]
Analyze the limitations of standard approaches to Turbulence when applied to edge cases in 2.005. Construct explicit counterexamples showing where intuitive reasoning fails.

\vspace{0.3cm}

\question[4]
Analyze the limitations of standard approaches to Thermodynamics when applied to edge cases in 2.005. Construct explicit counterexamples showing where intuitive reasoning fails.

\vspace{0.3cm}

\question[2]
Prove the fundamental theorem concerning Fluid Mechanics in the context of 2.005 (Thermal-Fluids Engineering I). State the theorem precisely, provide a complete proof, and discuss three important corollaries.

\vspace{0.3cm}

\question[5]
Analyze the limitations of standard approaches to Heat Transfer when applied to edge cases in 2.005. Construct explicit counterexamples showing where intuitive reasoning fails.

\vspace{0.3cm}

\question[4]
Analyze the limitations of standard approaches to Boundary Layers when applied to edge cases in 2.005. Construct explicit counterexamples showing where intuitive reasoning fails.

\vspace{0.3cm}

\question[3]
Construct an original proof-level problem involving Turbulence for 2.005 (Thermal-Fluids Engineering I) and provide a complete solution. Your problem should be at the level of a graduate qualifying exam.

\vspace{0.3cm}

\end{questions}

\newpage
\section*{Solutions}

\textbf{Solution 1:}

Solution approach: First recall the definitions and key theorems related to Thermodynamics from 2.005. Then apply the appropriate techniques to construct a rigorous argument. The solution must be self-contained and reference only the material from the course.

\vspace{0.5cm}

\textbf{Solution 2:}

The edge case analysis reveals the subtle assumptions underlying standard treatments of Fluid Mechanics. By constructing explicit counterexamples, we see that the usual hypotheses cannot be weakened without loss of the theorem's conclusion.

\vspace{0.5cm}

\textbf{Solution 3:}

The proof follows from the definitions and properties established in 2.005. The key insight is to apply the relevant theorem and verify the hypotheses hold. Detailed solution depends on the specific formulation of the theorem.

\vspace{0.5cm}

\textbf{Solution 4:}

The proof follows from the definitions and properties established in 2.005. The key insight is to apply the relevant theorem and verify the hypotheses hold. Detailed solution depends on the specific formulation of the theorem.

\vspace{0.5cm}

\textbf{Solution 5:}

The proof follows from the definitions and properties established in 2.005. The key insight is to apply the relevant theorem and verify the hypotheses hold. Detailed solution depends on the specific formulation of the theorem.

\vspace{0.5cm}

\textbf{Solution 6:}

The proof follows from the definitions and properties established in 2.005. The key insight is to apply the relevant theorem and verify the hypotheses hold. Detailed solution depends on the specific formulation of the theorem.

\vspace{0.5cm}

\textbf{Solution 7:}

Solution approach: First recall the definitions and key theorems related to Fluid Mechanics from 2.005. Then apply the appropriate techniques to construct a rigorous argument. The solution must be self-contained and reference only the material from the course.

\vspace{0.5cm}

\textbf{Solution 8:}

The edge case analysis reveals the subtle assumptions underlying standard treatments of Heat Transfer. By constructing explicit counterexamples, we see that the usual hypotheses cannot be weakened without loss of the theorem's conclusion.

\vspace{0.5cm}

\textbf{Solution 9:}

The edge case analysis reveals the subtle assumptions underlying standard treatments of Boundary Layers. By constructing explicit counterexamples, we see that the usual hypotheses cannot be weakened without loss of the theorem's conclusion.

\vspace{0.5cm}

\textbf{Solution 10:}

The edge case analysis reveals the subtle assumptions underlying standard treatments of Turbulence. By constructing explicit counterexamples, we see that the usual hypotheses cannot be weakened without loss of the theorem's conclusion.

\vspace{0.5cm}

\textbf{Solution 11:}

The proof follows from the definitions and properties established in 2.005. The key insight is to apply the relevant theorem and verify the hypotheses hold. Detailed solution depends on the specific formulation of the theorem.

\vspace{0.5cm}

\textbf{Solution 12:}

Solution approach: First recall the definitions and key theorems related to Fluid Mechanics from 2.005. Then apply the appropriate techniques to construct a rigorous argument. The solution must be self-contained and reference only the material from the course.

\vspace{0.5cm}

\textbf{Solution 13:}

The proof follows from the definitions and properties established in 2.005. The key insight is to apply the relevant theorem and verify the hypotheses hold. Detailed solution depends on the specific formulation of the theorem.

\vspace{0.5cm}

\textbf{Solution 14:}

The proof follows from the definitions and properties established in 2.005. The key insight is to apply the relevant theorem and verify the hypotheses hold. Detailed solution depends on the specific formulation of the theorem.

\vspace{0.5cm}

\textbf{Solution 15:}

The proof follows from the definitions and properties established in 2.005. The key insight is to apply the relevant theorem and verify the hypotheses hold. Detailed solution depends on the specific formulation of the theorem.

\vspace{0.5cm}

\textbf{Solution 16:}

Solution approach: First recall the definitions and key theorems related to Thermodynamics from 2.005. Then apply the appropriate techniques to construct a rigorous argument. The solution must be self-contained and reference only the material from the course.

\vspace{0.5cm}

\textbf{Solution 17:}

The edge case analysis reveals the subtle assumptions underlying standard treatments of Fluid Mechanics. By constructing explicit counterexamples, we see that the usual hypotheses cannot be weakened without loss of the theorem's conclusion.

\vspace{0.5cm}

\textbf{Solution 18:}

Solution approach: First recall the definitions and key theorems related to Heat Transfer from 2.005. Then apply the appropriate techniques to construct a rigorous argument. The solution must be self-contained and reference only the material from the course.

\vspace{0.5cm}

\textbf{Solution 19:}

The edge case analysis reveals the subtle assumptions underlying standard treatments of Boundary Layers. By constructing explicit counterexamples, we see that the usual hypotheses cannot be weakened without loss of the theorem's conclusion.

\vspace{0.5cm}

\textbf{Solution 20:}

The edge case analysis reveals the subtle assumptions underlying standard treatments of Turbulence. By constructing explicit counterexamples, we see that the usual hypotheses cannot be weakened without loss of the theorem's conclusion.

\vspace{0.5cm}

\textbf{Solution 21:}

The edge case analysis reveals the subtle assumptions underlying standard treatments of Thermodynamics. By constructing explicit counterexamples, we see that the usual hypotheses cannot be weakened without loss of the theorem's conclusion.

\vspace{0.5cm}

\textbf{Solution 22:}

The proof follows from the definitions and properties established in 2.005. The key insight is to apply the relevant theorem and verify the hypotheses hold. Detailed solution depends on the specific formulation of the theorem.

\vspace{0.5cm}

\textbf{Solution 23:}

The edge case analysis reveals the subtle assumptions underlying standard treatments of Heat Transfer. By constructing explicit counterexamples, we see that the usual hypotheses cannot be weakened without loss of the theorem's conclusion.

\vspace{0.5cm}

\textbf{Solution 24:}

The edge case analysis reveals the subtle assumptions underlying standard treatments of Boundary Layers. By constructing explicit counterexamples, we see that the usual hypotheses cannot be weakened without loss of the theorem's conclusion.

\vspace{0.5cm}

\textbf{Solution 25:}

Solution approach: First recall the definitions and key theorems related to Turbulence from 2.005. Then apply the appropriate techniques to construct a rigorous argument. The solution must be self-contained and reference only the material from the course.

\vspace{0.5cm}

\end{document}